\documentclass[12pt,letterpaper, onecolumn]{exam}
\usepackage{amsmath}
\usepackage{amssymb}
\usepackage{amsfonts}
\usepackage{graphicx}
\usepackage{subcaption}
\usepackage{float}
\usepackage[shortlabels]{enumitem}
\usepackage{hyperref}

\usepackage[lmargin=71pt, tmargin=1.2in]{geometry} %For centering solution box
\lhead{STAT 8561 - Homework 1\\}
\rhead{Tommy Stell\\}
% \chead{\hline} % Un-comment to draw line below header
\thispagestyle{empty}   %For removing header/footer from page 1

\begin{document}

\begingroup  
    \centering
    \LARGE Linear Statistical Analysis 1\\
    \LARGE Homework 1\\[0.5em]
    \large \today\\[0.5em]
    \large Tommy Stell\par
    
\endgroup
\rule{\textwidth}{0.4pt}
\pointsdroppedatright   %Self-explanatory
\printanswers
\renewcommand{\solutiontitle}{\noindent\textbf{Ans:}\enspace}   %Replace "Ans:" with starting keyword in solution box

\begin{questions}

    \question[8] Question 1\droppoints

    Consider the simple linear regression model
    \begin{equation*}
        Y_i = \beta_0 + \beta_1 x_i + \epsilon_i, \quad i=1,\ldots,5
    \end{equation*}

    with the following observed data
    \begin{equation*}
        \begin{array}{|c|c|c|}
           \hline i & x_i & y_i \\ \hline
            1 & 0 & 2 \\ \hline
            2 & 1 & 4 \\ \hline
            3 & 2 & 5 \\ \hline
            4 & 3 & 8 \\ \hline
            5 & 4 & 9 \\ \hline
        \end{array}
    \end{equation*}

    \begin{enumerate}[A.)]
        \item What is the design matrix?
        \item What is the dimension/rank of: 
        \begin{enumerate}[i.)]
            \item $\textbf{Y}$
            \item $\textbf{X}$
            \item $\beta$
            \item $\text{rank}(\textbf{X})$
        \end{enumerate}
        \item What is $\textbf{X}^\intercal\textbf{X}$ and $\textbf{X}^\intercal\textbf{Y}$?
        \item Assume $\mathbb{E}[\epsilon]=0$ and $\text{Var}(\epsilon)=\sigma^2I_5$. Evaluate the truthfulness of:
        \begin{enumerate}[i.)]
            \item $\mathbb{E}[\textbf{Y}]=\textbf{X}\beta$
            \item $\text{Var}(\textbf{Y})=\sigma^2I_5$
            \item $\mathbb{E}[\textbf{Y}]=0$
            \item The five errors all have variance $\sigma^2$
        \end{enumerate}
    \end{enumerate}
    
    \begin{solution}

    \begin{enumerate}[A.)]
        \item The design matrix is given by
        \begin{equation*}
            \textbf{X} = \left [\begin{array}{cc}
                1 & 0 \\
                1 & 1 \\
                1 & 2 \\
                1 & 3 \\
                1 & 4
            \end{array} \right]
        \end{equation*}
        \item \begin{enumerate}[i.)]
            \item $\textbf{Y}\in\mathbb{R}^{5\times 1}$ with rank $=1$
            \item $\textbf{X}\in\mathbb{R}^{5\times 2}$ with rank $=2$
            \item $\beta\in\mathbb{R}^{2\times 1}$ with rank $=1$
            \item $\text{rank}(\textbf{X})\in\mathbb{R}^1$ and is a scalar
        \end{enumerate}
        \item 
        \begin{equation*}
            \textbf{X}^\intercal\textbf{X} = \left [\begin{array}{ccccc}
                1 & 1 &  1 & 1 & 1 \\ 0 & 1 & 2 & 3 & 4
            \end{array} \right] \left [\begin{array}{cc}
                1 & 0 \\
                1 & 1 \\
                1 & 2 \\
                1 & 3 \\
                1 & 4
            \end{array} \right] = \left [\begin{array}{cc}
                5 & 10 \\
                10 & 30 
            \end{array} \right]
        \end{equation*}
        \begin{equation*}
            \textbf{X}^\intercal\textbf{Y} = \left [\begin{array}{ccccc}
                1 & 1 &  1 & 1 & 1 \\ 0 & 1 & 2 & 3 & 4
            \end{array} \right]  \left [\begin{array}{c}
                2  \\
                4  \\
                5  \\
                8  \\
                9 
            \end{array} \right] = \left [\begin{array}{c}
                28 \\ 74
            \end{array} \right]
        \end{equation*}
        \item \begin{enumerate}[i.)]
            \item True
            
            Since $\textbf{Y}=\textbf{X}\beta + \epsilon$, we can take the expectation of both sides resulting in
            \begin{align*}
                \mathbb{E}[\textbf{Y}] & = \mathbb{E}[\textbf{X}\beta + \epsilon] \\
                \mathbb{E}[\textbf{Y}] & = \mathbb{E}[\textbf{X}\beta] +\mathbb{E}[ \epsilon] \\
                \mathbb{E}[\textbf{Y}] & = \mathbb{E}[\textbf{X}\beta] + 0 \\
                \mathbb{E}[\textbf{Y}] & = \textbf{X}\beta \quad \text{since} \quad \mathbb{E}[\textbf{X}\beta] = \textbf{X}\beta \\
            \end{align*}
            \item True
            \begin{align*}
                \text{Var}(\textbf{Y}) &= \text{Var}(\textbf{X}\beta + \epsilon) \\
                \text{Var}(\textbf{Y}) &= \text{Var}( \epsilon) \quad \text{since} \quad \text{Var}(\textbf{X}\beta)=0\\
                \text{Var}(\textbf{Y}) &= \sigma^2I_5 \quad \text{since} \quad \text{Var}( \epsilon) = \sigma^2I_5
            \end{align*}
            \item Usually False
            
            Since we showed that $\mathbb{E}[\textbf{Y}] = \textbf{X}\beta$, this is only true in the case that $\textbf{X}\beta = 0$ which comes from either the model predicting a constant output at 0 or if the model has only weights of 0.
            \item True
            
            Let $\epsilon$ have the form
            \begin{equation*}
                \epsilon = \left[\begin{array}{c}
                \epsilon_1  \\
                \epsilon_2  \\
                \epsilon_3  \\
                \epsilon_4  \\
                \epsilon_5 
            \end{array} \right]
            \end{equation*}
            then, we know that the $\text{Var}(\epsilon)=\sigma^2I_5$ or,
            \begin{equation*}
               \text{Var}(\epsilon)= \text{Var}\left(\left[\begin{array}{c}
                \epsilon_1  \\
                \epsilon_2  \\
                \epsilon_3  \\
                \epsilon_4  \\
                \epsilon_5 
            \end{array} \right] \right) = \left [\begin{array}{ccccc}
                \sigma^2 & 0 & 0 & 0 & 0 \\ 0 & \sigma^2 & 0 & 0 & 0 \\
                0 & 0 & \sigma^2 & 0 & 0 \\ 0 & 0 & 0 & \sigma^2 & 0 \\
                0 & 0 & 0 & 0 & \sigma^2
            \end{array} \right]
            \end{equation*}
            so each error has the same variance of exactly $\sigma^2$.
            \end{enumerate}
    \end{enumerate}

        $\square$
    \end{solution}

    \newpage
    \question[12] Question 2\droppoints

    Continue using the data and design matrix from question 1.

    \begin{enumerate}[A.)]
        \item Starting from \\ \\
        $S(\beta)=(\textbf{Y}-\textbf{X}\beta)^\intercal(\textbf{Y}-\textbf{X}\beta)$, \\ \\
        derive the normal equations \\ \\
        $\textbf{X}^\intercal \textbf{X}\hat\beta = \textbf{X}^\intercal \textbf{Y}$ \\ \\
        and then state the OLS when $\textbf{X}$ has full column rank.
        \item Using the quantities from question 1, compute \\ \\
        $\hat\beta = \left[ \begin{array}{c}
             \hat\beta_0  \\ \hat\beta_1
        \end{array}\right]$ \\ \\
        Identify the values of $\hat\beta_0$ and $\hat\beta_1$, then write out the estimator $\hat Y$.
        \item For the least squares residual vector, $\textbf{e}=\textbf{Y}-\hat{\textbf{Y}}$, evaluate the truthfulness of each expression and justify geometrically: \begin{enumerate}[i.)]
            \item $\textbf{X}\textbf{e}=0$
            \item $\textbf{X}^\intercal \textbf{e}=0$
            \item $\textbf{e}^\intercal \textbf{e}=0$
            \item $\textbf{H}\textbf{e}=\textbf{e}$
        \end{enumerate}
    \end{enumerate}
    
    \begin{solution}

        \begin{enumerate}[a.)]
            \item Given $S(\beta)=(\textbf{Y}-\textbf{X}\beta)^\intercal(\textbf{Y}-\textbf{X}\beta)$ then by expansion,
            \begin{equation*}
                S(\beta) = \textbf{Y}^\intercal \textbf{Y} - \textbf{Y}^\intercal \textbf{X}\beta - \beta^\intercal \textbf{X}^\intercal\textbf{Y} + \beta^\intercal \textbf{X}^\intercal \textbf{X}\beta
            \end{equation*}

            The term $\textbf{Y}^\intercal\textbf{X}\beta$ has terms $\textbf{Y}^\intercal\in\mathbb{R}^{1\times n}$, $\textbf{X}\in\mathbb{R}^{n\times p}$, and $\beta\in\mathbb{R}^{p\times1}$ which results in the product being in $\mathbb{R}^{1\times1}\subseteq\mathbb{R}$. Thus, 
            \begin{equation*}
                \textbf{Y}^\intercal\textbf{X}\beta=(\textbf{Y}^\intercal\textbf{X}\beta)^\intercal = \beta^\intercal\textbf{X}^\intercal \textbf{Y}
            \end{equation*}

            Together, this simplifies $S(\beta)$ into,
            \begin{equation*}
                S(\beta)= \textbf{Y}^\intercal\textbf{Y} - 2 \beta^\intercal\textbf{X}^\intercal\textbf{Y} + \beta^\intercal\textbf{X}^\intercal\textbf{X}\beta
            \end{equation*}

            Since we are seeking the minimum value of $S(\beta)$ given $\beta$, we can take the partial derivative of $S(\beta)$ with respect to $\beta$. We can use the identity that was blatantly stolen from \href{https://en.wikipedia.org/wiki/Matrix_calculus}{Wikipedia},
            \begin{equation*}
                \frac{\partial (\textbf{Ax}+\textbf{b})^\intercal\textbf{C}(\textbf{Dx} + \textbf{e})}{\partial \textbf{x}} = \textbf{D}^\intercal\textbf{C}^\intercal(\textbf{Ax}+\textbf{b})+\textbf{A}^\intercal \textbf{C}(\textbf{Dx}+\textbf{e})
            \end{equation*}

            with $\textbf{b}=\textbf{e}=0$, $\textbf{A}=\textbf{D}=\textbf{I}$, and $\textbf{C}=\textbf{X}^\intercal\textbf{X}$, we get, 

            \begin{equation*}
                \frac{\partial S(\beta)}{\partial\beta} = - 2 \textbf{X}^\intercal\textbf{Y} + 2\textbf{X}^\intercal\textbf{X}\beta
            \end{equation*}

            The minimum, is of course when the derivative is equal to 0, and thus,
            \begin{equation*}
                - 2 \textbf{X}^\intercal\textbf{Y} + 2\textbf{X}^\intercal\textbf{X}\hat{\beta} = 0
            \end{equation*}

            which leaves the final normal equations of,
            \begin{equation*}
                \textbf{X}^\intercal\textbf{Y} = \textbf{X}^\intercal\textbf{X}\hat{\beta}
            \end{equation*}

            The OLS estimator is thus,
            \begin{equation*}
                \hat{\beta} = ( \textbf{X}^\intercal\textbf{X})^{-1}\textbf{X}^\intercal\textbf{Y}
            \end{equation*}

        \item To find the value of $\hat{\beta}$, we can set $\hat{\beta}= ( \textbf{X}^\intercal\textbf{X})^{-1}\textbf{X}^\intercal\textbf{Y}$. Copying the above from question 1, we have,
        \begin{align*}
            \textbf{X}^\intercal\textbf{X} & = \left [\begin{array}{cc}
                5 & 10 \\
                10 & 30 
            \end{array} \right] \\
            \textbf{X}^\intercal\textbf{Y} & = \left [\begin{array}{c}
                28 \\ 74
            \end{array} \right]
        \end{align*}

        Thus, it is easy to compute the inverse of $\textbf{X}^\intercal\textbf{X}$, assuming it exists, and can be done as,
        \begin{equation*}
            (\textbf{X}^\intercal\textbf{X})^{-1}=\frac{1}{150 - 100} \left [\begin{array}{cc}
                30 & -10 \\
                -10 & 5 
            \end{array} \right] =  \left [\begin{array}{cc}
                \frac{3}{5} & -\frac{1}{5} \\
                -\frac{1}{5} & \frac{1}{10} 
            \end{array} \right]
        \end{equation*}

        Thus, $\hat{\beta}$ is given by,
        \begin{align*}
            \hat{\beta} & = ( \textbf{X}^\intercal\textbf{X})^{-1}\textbf{X}^\intercal\textbf{Y} \\
            & = \left [\begin{array}{cc}
                \frac{3}{5} & -\frac{1}{5} \\
                -\frac{1}{5} & \frac{1}{10} 
            \end{array} \right]\left [\begin{array}{c}
                28 \\ 74
            \end{array} \right] \\ 
            & = \left [\begin{array}{c}
                2 \\ \frac{9}{5}
            \end{array} \right]
        \end{align*}

        Using the found value of $\hat{\beta}$, we can say that the estimator is thus,
        \begin{align*}
            \hat{\textbf{Y}} &= \textbf{X}\hat{\beta} \\ 
             \hat{\textbf{Y}}&= \left [\begin{array}{cc}
                1 & 0 \\
                1 & 1 \\
                1 & 2 \\
                1 & 3 \\
                1 & 4
            \end{array} \right]\left [\begin{array}{c}
                2 \\ \frac{9}{5}
            \end{array} \right] \\
             \hat{\textbf{Y}} & = \left[\begin{array}{c}
                2  \\
                3.8  \\
                5.6  \\
                7.4  \\
                9.2 
            \end{array} \right]
        \end{align*}

        The final estimator is thus,
        \begin{equation*}
            \hat{\textbf{Y}} = 2 + \frac{9}{5}\textbf{X}
        \end{equation*}

        \item The error vector $\textbf{e}$ is thus,
        \begin{align*}
            \textbf{e}&=\textbf{Y}- \hat{\textbf{Y}} \\
            \textbf{e} & = \left[\begin{array}{c}
                2  \\
                4  \\
                5  \\
                8  \\
                9 
            \end{array} \right] - \left[\begin{array}{c}
                2  \\
                3.8  \\
                5.6  \\
                7.4  \\
                9.2 
            \end{array} \right] = \left[\begin{array}{c}
                0  \\
                0.2  \\
                -0.6  \\
                0.6  \\
                -0.2 
            \end{array} \right]
        \end{align*}

        \begin{enumerate}[i.)]
            \item False 

            The product $\textbf{Xe}$ is meaningless as $\textbf{X}\in\mathbb{R}^{n\times 2}$ and $\textbf{e}\in\mathbb{R}^{n\times 1}$. They do not have a meaningful product via usual matrix multiplication.

            \item True

            This will at least have a product which exists. Computing this, we get,
            \begin{equation*}
                \textbf{X}^\intercal\textbf{e} = \left [\begin{array}{ccccc}
                1 & 1 & 1 & 1 & 1 \\
                0 & 1 & 2 & 3 & 4
            \end{array} \right]\left[\begin{array}{c}
                0  \\
                0.2  \\
                -0.6  \\
                0.6  \\
                -0.2 
            \end{array} \right] = \left[\begin{array}{c}
                 0  \\  0                   
            \end{array} \right]
            \end{equation*}

            Geometrically, the fitted model is the best that can be achieved. There is still some residual and the amount left is not inside the span of $\textbf{X}$. The residual is orthogonal to the span of $\textbf{X}$ thus leaving the inner product 0.

            \item False

            This is the magnitude of the residual vector. This can be equal to 0 but is not usually. It is only equal to 0 if the actual data is collinear. 

            \item False

            With $\textbf{H}$ defined as,
            \begin{equation*}
                \textbf{H}=\textbf{X}(\textbf{X}^\intercal\textbf{X})^{-1}\textbf{X}^\intercal
            \end{equation*}

            We can right multiply by $\textbf{e}$ to get,
            \begin{align*}
                \textbf{He}& = \textbf{X}(\textbf{X}^\intercal\textbf{X})^{-1}\textbf{X}^\intercal\textbf{e} \\
                \textbf{He}& = \textbf{X}(\textbf{X}^\intercal\textbf{X})^{-1}(\textbf{X}^\intercal\textbf{e}) \\
                \textbf{He}& = \textbf{X}(\textbf{X}^\intercal\textbf{X})^{-1}0 \\ \textbf{He}&=0
            \end{align*}

            Since $\textbf{H}$ depends only on $\textbf{X}$, and with $\textbf{e}$ outside the span of $\textbf{X}$, how could $\textbf{H}$ ever reach a point that is orthogonal to it? It can't leaving the identity as false.
            
        \end{enumerate}
            
        \end{enumerate}
    
        $\square$
        
    \end{solution}

    \newpage

    \question[12] Question 3\droppoints

    Continue using the data and fitted regression model from before.

    \begin{enumerate}[a.)]
        \item Compute both $\hat{\textbf{Y}}=\textbf{X}\hat{\beta}$ and $\textbf{e}=\textbf{Y}-\hat{\textbf{
       Y 
        }}$.
        \item Determine the truth value of:
        \begin{enumerate}[i.)]
            \item The residuals satisfy
            \begin{equation*}
                \sum\limits_{i=1}^{n}\textbf{e}_i=0
            \end{equation*}
            \item The fitted values satisfy $\overline{\hat{\textbf{Y}}}=\overline{\textbf{Y}}$.
            \item The vector composing of all 1's (i.e. $\textbf{1}_n$) belongs to $\mathcal{C}(\textbf{X}).$
            \item These properties necessarily hold for every regression model without an intercept.
        \end{enumerate}
        \item Define $\textbf{M}=\textbf{I}_n - \textbf{H}$. 
        Determine the truth value of:
        \begin{enumerate}[i.)]
            \item $\textbf{e}=\textbf{M}\textbf{Y}$
            \item $\textbf{M}^\intercal = \textbf{M}$
            \item $\textbf{M}^2 = \textbf{M}$
            \item $\textbf{HM}=\textbf{0}$
        \end{enumerate}
        \item Suppose now that a design matrix $\textbf{X}$ does not have full column rank.
        Determine the truth value of:
        \begin{enumerate}[i.)]
            \item $(\textbf{X}^\intercal\textbf{X})^{-1}$ necessarily exists.
            \item The least squares coefficient vector $\hat{\beta}$ may not be unique. 
            \item The fitted vector $\hat{\textbf{Y}}$ is still unique.
            \item The residual vector $\textbf{e}$ is still unique.
        \end{enumerate}
    \end{enumerate}

    \begin{solution}

    \begin{enumerate}[a.)]
        \item Copying from above,

        \begin{align*}
            \hat{\textbf{Y}} &= \textbf{X}\hat{\beta} \\ 
             \hat{\textbf{Y}}&= \left [\begin{array}{cc}
                1 & 0 \\
                1 & 1 \\
                1 & 2 \\
                1 & 3 \\
                1 & 4
            \end{array} \right]\left [\begin{array}{c}
                2 \\ \frac{9}{5}
            \end{array} \right] \\
             \hat{\textbf{Y}} & = \left[\begin{array}{c}
                2  \\
                3.8  \\
                5.6  \\
                7.4  \\
                9.2 
            \end{array} \right]
        \end{align*}

        and,

        \begin{align*}
            \textbf{e}&=\textbf{Y}- \hat{\textbf{Y}} \\
            \textbf{e} & = \left[\begin{array}{c}
                2  \\
                4  \\
                5  \\
                8  \\
                9 
            \end{array} \right] - \left[\begin{array}{c}
                2  \\
                3.8  \\
                5.6  \\
                7.4  \\
                9.2 
            \end{array} \right] = \left[\begin{array}{c}
                0  \\
                0.2  \\
                -0.6  \\
                0.6  \\
                -0.2 
            \end{array} \right]
        \end{align*}

        \item 
        \begin{enumerate}[i.)]
            \item True

            Taking $\textbf{e}$, we get,
            \begin{equation*}
                \left[\begin{array}{c}
                0  \\
                0.2  \\
                -0.6  \\
                0.6  \\
                -0.2 
            \end{array} \right] \implies 0 + 0.2 -0.6 + 0.6 - 0.2 = 0
            \end{equation*}
            
            \item True

            \begin{align*}
                \overline{\textbf{Y}} & = \frac{2 + 4 + 5 +8 + 9}{5} = 5.6 \\
                \overline{\hat{\textbf{Y}}} & = \frac{2 + 3.8 + 5.6 + 7.4 + 9.2}{5} = 5.6
            \end{align*}

            \item True

            If we partition $\textbf{X}$ as $\textbf{X}=[ \textbf{x}_1 \quad \textbf{x}_2]$, then $\textbf{x}_1 = \textbf{1}_5$.

            \item False

            The identity of,
            \begin{equation*}
                \sum\limits_{i=1}^{n}\textbf{e}_i=0
            \end{equation*}
            does not hold if the intercept drops. 
        \end{enumerate}
        \item Here, $\textbf{M} = \textbf{I}_5 - \textbf{H}$,
        \begin{align*}
            \textbf{M} &= \left[\begin{array}{ccccc}
                1 & 0 & 0 & 0 & 0 \\
                0 & 1 & 0 & 0 & 0 \\
                0 & 0 & 1 & 0 & 0 \\
                0 & 0 & 0 & 1 & 0 \\
                0 & 0 & 0 & 0 & 1
            \end{array} \right] - \left[\begin{array}{ccccc}
                 0.6 & 0.4 & 0.2 & 0.0 & -0.2 \\
                 0.4 & 0.3 & 0.2 & 0.1 &  0.0 \\
                 0.2 & 0.2 & 0.2 & 0.2 &  0.2 \\
                 0.0 & 0.1 & 0.2 & 0.3 &  0.4 \\
                -0.2 & 0.0 & 0.2 & 0.4 &  0.6
            \end{array}\right] \\ \\
            & = \left[\begin{array}{ccccc}
                 0.4 & -0.4 & -0.2 &  0.0 &  0.2 \\
                -0.4 &  0.7 & -0.2 & -0.1 &  0.0 \\
                -0.2 & -0.2 &  0.8 & -0.2 & -0.2 \\
                 0.0 & -0.1 & -0.2 &  0.7 & -0.4 \\
                 0.2 &  0.0 & -0.2 & -0.4 &  0.4
            \end{array}\right]
        \end{align*}
        \begin{enumerate}[i.)]
            \item True

            Calculating directly, 
            \begin{align*}
                \textbf{MY} & = \left[\begin{array}{ccccc}
                 0.4 & -0.4 & -0.2 &  0.0 &  0.2 \\
                -0.4 &  0.7 & -0.2 & -0.1 &  0.0 \\
                -0.2 & -0.2 &  0.8 & -0.2 & -0.2 \\
                 0.0 & -0.1 & -0.2 &  0.7 & -0.4 \\
                 0.2 &  0.0 & -0.2 & -0.4 &  0.4
            \end{array}\right]\left[\begin{array}{c}
                2  \\
                4  \\
                5  \\
                8  \\
                9 
            \end{array} \right] \\
            & = \left[\begin{array}{c}
                0  \\
                0.2  \\
                -0.6  \\
                0.6  \\
                -0.2 
            \end{array} \right] \\
            & = \textbf{e}
            \end{align*}
            \item True

            Since $\textbf{M}$ is defined as $\textbf{M}=\textbf{I} - \textbf{H}$, we can use the properties of both $\textbf{I}$ and $\textbf{H}$.
            \begin{align*}
                \textbf{M}^\intercal & = (\textbf{I} - \textbf{H})^\intercal \\
                &= \textbf{I}^\intercal - \textbf{H}^\intercal \\
                & = \textbf{I} - \textbf{H} \quad \text{since} \quad \textbf{I}^\intercal = \textbf{I} \quad \& \quad \textbf{H}^\intercal = \textbf{H} \\
                & = \textbf{M}
            \end{align*}

            \item True

            \begin{align*}
                \textbf{M}^2 &= (\textbf{I} - \textbf{H})(\textbf{I} - \textbf{H}) \\
                & = \textbf{I}^2 - \textbf{HI} - \textbf{IH} + \textbf{H}^2 \\
                & = \textbf{I} - 2\textbf{H} + \textbf{H} \quad \text{since} \quad \textbf{H}^2=\textbf{H} \\
                &= \textbf{I} - \textbf{H} \\
                &= \textbf{M}
            \end{align*}
            \item True

            \begin{align*}
                \textbf{HM} & = \textbf{H}(\textbf{I} - \textbf{H}) \\
                &= \textbf{HI} - \textbf{H}^2 \\
                &= \textbf{H} - \textbf{H} \quad \text{since} \quad \textbf{H}^2=\textbf{H} \\
                & = 0
            \end{align*}
        \end{enumerate}
        \item We assume that $\textbf{X}$ does not have full column rank.
        \begin{enumerate}[i.)]
            \item False

            This is due to the following properties. Since $\textbf{X}$ does not have full column rank, its determinant is 0. The value $(\textbf{X}^\intercal\textbf{X})^{-1}$ exists iff that $\det[\textbf{X}^\intercal\textbf{X}]\neq0$.
            
            \begin{align*}
               & \text{rank}(\textbf{X}^\intercal\textbf{X}) =  \text{rank}(\textbf{X}) \\
               & \text{rank}(\textbf{X}) < p \implies \text{rank}(\textbf{X}^\intercal\textbf{X}) < p \\
               & \text{rank}(\textbf{X}^\intercal\textbf{X}) < p\implies  \det(\textbf{X}^\intercal\textbf{X})=0
            \end{align*}

            Which is a contradiction leading to False.

            \item True

            The values of $\beta$ do not need to be unique.
            
            \item True

            Due to $\textbf{X}$ not having full column rank, the $\hat{\beta}$ that is chosen, while not unique, does point toward a unique $\hat{\textbf{Y}}$ which minimizes the least squares problem.

            \item True

            Taking $\hat{\textbf{Y}}$ to be unique and $\textbf{Y}$ to be constant, then necessarily, $\textbf{e}=\textbf{Y}-\hat{\textbf{Y}}$ is also unique since it is just an affine shift of $\hat{\textbf{Y}}$.
        \end{enumerate}
    \end{enumerate}

        
        $\square$
        
    \end{solution}

    \newpage

    \question[8] Question 4\droppoints

    Using the built-in model in R, $\textit{cars}$, consider the model,
    \begin{equation*}
        \text{dist}_i=\beta_0 + \beta_1\text{speed}_i + \epsilon_i
    \end{equation*}

    \begin{enumerate}[a.)]
        \item Fit the model using linear regression. Identify the vector $\beta$ and interpret $\beta_1$.
        \item Construct $\textbf{Y}$ and $\textbf{X}$ manually. Reproduce the OLS coefficients using matrix operations.
    \end{enumerate}

    \begin{solution}

        \begin{enumerate}[a.)]
            \item 
        
            Running the code (in accompanying file), I get the $\beta$ vector to be,
            \begin{equation*}
                \beta = \left[\begin{array}{c}
                    -17.579 \\ 3.932 
                \end{array} \right]
            \end{equation*}
    
            This gives an overall linear regression of,
            \begin{equation*}
                \widehat{\text{dist}}_i=-17.579 + 3.932\text{speed}_i
            \end{equation*}
    
            What $\beta_1$ means is that the distance required to stop goes up by 3.932 feet for every mph increase of speed. 
    
            \item 

            The data and targets are shown,

            \begin{equation*}
                \textbf{X} = \left[\begin{array}{cc}
                    1 & 4 \\
                    1 & 4 \\
                    1 & 7 \\
                    \vdots & \vdots \\
                    1 & 25
                \end{array} \right] \quad \& \quad \textbf{Y} = \left[\begin{array}{c}
                    2 \\ 10 \\ 4 \\
                    \vdots  \\
                    85
                \end{array} \right]
            \end{equation*}

            For brevity, I am not going to write those out explicitly but I am going to use them to hand calculate the $\beta$ vector. We know that $\hat \beta$ is given by,
            \begin{equation*}
                \hat\beta = ( \textbf{X}^\intercal\textbf{X})^{-1}\textbf{X}^\intercal\textbf{Y}
            \end{equation*}

            We can painfully do the matrix operations to get that,
            \begin{equation*}
                \textbf{X}^\intercal\textbf{X} = \left[\begin{array}{cc}
                    50 & 770 \\ 770 & 13228
                \end{array} \right] \quad \& \quad \textbf{X}^\intercal\textbf{Y} = \left[\begin{array}{c}
                    2149 \\ 38482
                \end{array} \right]
            \end{equation*}

            These are easier to work with. By taking the inverse we get,
            \begin{equation*}
                ( \textbf{X}^\intercal\textbf{X})^{-1} = \frac{1}{34250} \left[\begin{array}{cc}
                    6614 & -385 \\ -385 & 25
                \end{array} \right]
            \end{equation*}

            We thus get the multiplication problem of,
            \begin{align*}
                \hat\beta &= ( \textbf{X}^\intercal\textbf{X})^{-1}\textbf{X}^\intercal\textbf{Y}\\&= \frac{1}{34250} \left[\begin{array}{cc}
                    6614 & -385 \\ -385 & 25
                \end{array} \right]\left[\begin{array}{c}
                    2149 \\ 38482
                \end{array} \right]
                \\ &=\left[\begin{array}{c}
                    -17.579 \\ 3.932
                \end{array} \right]
            \end{align*}

            Which is exactly what our R code produced.

        \end{enumerate}        
        
        $\square$
        
    \end{solution}

    
\end{questions}
\end{document}
